Con el objetivo de facilitar la comprension del sistema y organizar las funcionalidades de forma coherente, los requisitos se han agrupado en modulos funcionales que representan las principales areas de operacion del sistema. Estos modulos se describen a continuacion de forma general, mientras que en los apartados posteriores se detallan sus requisitos funcionales y de informacion asociados.

La descomposicion del sistema en modulos funcionales responde a los principios de diseno arquitectonico en ingenieria del software, que promueven la organizacion del sistema en componentes con responsabilidades bien definidas, alta cohesion y bajo acoplamiento, con el objetivo de facilitar su comprension, mantenimiento y evolucion \citep[Ch.~6--7]{sommerville}.

\subsection{M1. Gestion de acceso, alta y autorizacion de usuarios}

Este modulo agrupa las funcionalidades relacionadas con la identificacion de los usuarios del sistema, el alta de nuevas cuentas, la autenticacion y el control de acceso segun el perfil asignado. Su finalidad es garantizar que unicamente usuarios autorizados puedan utilizar la aplicacion y que cada uno acceda exclusivamente a las funciones que le correspondan dentro del sistema.

Asimismo, este modulo contempla la gestion de roles, el tratamiento de solicitudes de registro, la administracion del estado de las cuentas y el registro de eventos de acceso y acciones administrativas, con el fin de aportar seguridad, trazabilidad y control sobre el uso de la plataforma.

Este modulo constituye la base del sistema, ya que establece los mecanismos de acceso y control sobre los que se apoyan el resto de funcionalidades. Su correcta definicion permite garantizar la seguridad del sistema, la coherencia en la gestion de usuarios y la trazabilidad de las operaciones realizadas.

\subsection{M2. Gestion comercial y clientes}

Este modulo agrupa las funcionalidades relacionadas con la gestion de los clientes profesionales (peluquerias) y con la operativa diaria del equipo comercial. Su objetivo es centralizar la informacion relevante de cada cliente, mantener controlada su asignacion a un comercial concreto y facilitar la planificacion del trabajo de campo mediante visitas, rutas y estimaciones operativas de reparto.

Ademas de la informacion basica y comercial de cada cliente, este modulo incorpora la geolocalizacion de los establecimientos, la definicion de franjas horarias validas para visitas, la configuracion operativa de la jornada de cada comercial y la generacion dinamica de una planificacion diaria. De este modo, el sistema no actua unicamente como repositorio de fichas de cliente, sino tambien como herramienta de apoyo a la toma de decisiones durante la jornada, calculando orden de visitas, horas estimadas de llegada, margenes disponibles y situaciones que requieren replanificacion.

\subsection{M3. Catalogo, productos y coloracion}

Este modulo agrupa las funcionalidades relacionadas con la consulta y gestion del catalogo de productos profesionales comercializados por el distribuidor. Su objetivo es proporcionar a los clientes profesionales y a los agentes comerciales una forma estructurada de acceder a la informacion de los productos, sus caracteristicas tecnicas y su clasificacion dentro de las distintas lineas comerciales.

Asimismo, el modulo contempla la organizacion jerarquica del catalogo mediante categorias y lineas comerciales, asi como la gestion de recursos de apoyo, documentacion tecnica y materiales complementarios asociados a productos o gamas. Tambien incluye herramientas especificas para la consulta de cartas de color y referencias cromaticas utilizadas en los procesos de coloracion capilar, lo que refuerza su utilidad tanto comercial como tecnica.

\subsection{M4. Pedidos, entregas y cobros}

Este modulo agrupa las funcionalidades relacionadas con la gestion del proceso comercial de venta de productos, desde la creacion de borradores de pedido hasta la entrega fisica y el seguimiento minimo del cobro asociado. Su objetivo es permitir que clientes profesionales, comerciales y administradores trabajen sobre un mismo flujo trazable de pedido, con visibilidad del estado del encargo, de su reparto y de su situacion de cobro.

En la implementacion actual, el modulo se apoya en pedidos con lineas de producto, estados de ciclo de vida, vinculacion opcional a visitas de reparto y un seguimiento minimo del cobro mediante estados \texttt{pending}/\texttt{paid}, metodo, fecha y notas. La fase logistica no se modela como un subsistema aislado de entregas, sino integrada en las visitas comerciales de tipo reparto, con validacion de cierre mediante QR y con ETA visible para cliente solo cuando existe un reparto real y coherente. Quedan explicitamente fuera de este alcance minimo los pagos parciales, la facturacion y la conciliacion avanzada.

\subsection{M5. Gestion tecnica del salon y seguimiento de servicios}

Este modulo agrupa las funcionalidades destinadas a registrar y consultar informacion tecnica sobre los servicios realizados en las peluquerias clientes del sistema. Su objetivo es permitir a los profesionales del salon llevar un seguimiento estructurado de los trabajos realizados a sus clientes, especialmente aquellos relacionados con productos de la marca, facilitando la consulta de historiales tecnicos y el analisis del uso de productos.

En este sentido, el modulo contempla la gestion de clientes finales del salon, los servicios o tratamientos realizados, sus fichas tecnicas asociadas y los productos empleados en cada intervencion. Ademas, incorpora soporte para la generacion de sugerencias de productos y para la elaboracion de comunicaciones personalizadas a partir de la informacion tecnica registrada, reforzando asi tanto la continuidad del trabajo profesional como la fidelizacion del cliente final.

\subsection{M6. Promociones, notificaciones y formaciones}

Este modulo agrupa las funcionalidades relacionadas con la comunicacion entre el distribuidor y las peluquerias usuarias del sistema. Su objetivo es facilitar la difusion de promociones comerciales, la organizacion de actividades formativas y el envio de notificaciones o recordatorios relevantes para los usuarios de la aplicacion.

Ademas, este modulo permite segmentar clientes mediante categorias comerciales, asociar promociones a productos, lineas o grupos de clientes, y gestionar avisos vinculados a eventos relevantes del sistema. De este modo, no solo cubre la comunicacion promocional, sino tambien la comunicacion operativa y la gestion de acciones formativas e inscripciones.

\subsection{M7. Administracion e integraciones}

Este modulo agrupa las funcionalidades destinadas a la administracion general del sistema y a la integracion con herramientas o servicios externos utilizados en la actividad del distribuidor. Su objetivo es permitir la configuracion y mantenimiento de la aplicacion, facilitar el analisis de la actividad comercial y permitir la interoperabilidad con sistemas externos utilizados en la operativa del negocio.

En particular, este modulo contempla la gestion de parametros generales de configuracion, el registro y seguimiento de integraciones externas y las operaciones de exportacion o sincronizacion de informacion. Asimismo, incorpora soporte a la generacion de propuestas de pedido a proveedor a partir de la informacion registrada en la plataforma o introducida manualmente, alineandose con la operativa habitual del distribuidor.

\subsection{M8. Funcionalidades adicionales}

Este modulo recoge un conjunto de funcionalidades identificadas durante el proceso de analisis de requisitos que por ellas solas no forman parte de ningun modulo principal. Estas funcionalidades buscan ampliar el valor de la aplicacion para los clientes mediante el uso de tecnologias avanzadas y nuevas herramientas de apoyo al trabajo del peluquero.

Entre estas funcionalidades se incluyen, de forma potencial, herramientas de recomendacion de tratamientos, simulacion visual de coloracion, reconocimiento de productos mediante camara y gestion de reservas online por parte de clientes finales. Estas funcionalidades pueden llegar a entrar, o no, en el alcance del proyecto dependiendo de la disponibilidad temporal del mismo. Si no llegasen a incluirse en la primera version del sistema, podrian ser consideradas para futuras actualizaciones o mejoras de la aplicacion.
